\label{sec:mip_model}

In this section, we present the Mixed-Integer Programming (MIP) formulation for the fixture layout optimization problem. Unlike the CP model, the MIP formulation does not rely on global constraints. Instead, it represents the problem as a system of linear constraints over continuous and binary decision variables. This
relaxation comes from the nature of MIP solving which, unlike most CP solvers, naturally handles
non-negative continuous variables.


\subsection{Variables}
Variables \(x_i\) and \(y_i\), for \(i=1,\dots,\nmax\), have the same semantics as those defined in Section~\ref{sec:cp_vars}, with the only difference that they are continuous rather than discrete.

Fixture types are encoded by a binary matrix \(T\) of size \(\nmax \times \ntypes\), where

\[
T_{i,j} =
\begin{cases}
	1 & \text{if fixture \(i\) is selected and has type \(j\),}\\
	0 & \text{otherwise.}
\end{cases}
\]

To represent the width, height, and type of each fixture, auxiliary variables \(w_i\), \(h_i\), and \(t_i\) are introduced for each \(i \in [1 \ldots \nmax]\).

To linearize the constraints presented in Section~\ref{sec:cp_cons}, the following auxiliary binary variables are defined:

\begin{itemize}
	\item \(\used_i\), which indicates whether fixture \(i\) is selected;
	\item \(\samebar_{i,j}\), which indicates whether fixtures \(i\) and \(j\) are placed on the same bar;
	\item \(\aboveinbar_{i,j}\) and \(\belowinbar_{i,j}\), which indicate whether fixture \(i\) is positioned above or below fixture \(j\) on the same bar;
	\item \(\nooverlapabove_{i,j}\), \(\nooverlapbelow_{i,j}\), \(\nooverlapleft_{i,j}\), and \(\nooverlapright_{i,j}\), which indicate whether fixture \(i\) is positioned above, below, to the left, or to the right of fixture \(j\), respectively;
	\item \(\aboveext_{i,j}\), \(\belowext_{i,j}\), \(\leftext_{i,j}\), and \(\rightext_{i,j}\), which indicate whether fixture \(i\) is positioned above, below, to the left, or to the right of extrusion area \(j\), respectively.
\end{itemize}


\subsection{Constraints}
For each fixture \(i \in [1..\nmax]\), the following equations are imposed:

\[
\used_i = \sum_{j=1}^{\ntypes} T_{i,j}, \qquad
w_i = \sum_{j=1}^{\ntypes} \wfix_j \cdot T_{i,j}, \qquad
h_i = \sum_{j=1}^{\ntypes} \hfix_j \cdot T_{i,j}, \qquad
t_i = \sum_{j=1}^{\ntypes} j \cdot T_{i,j}.
\]

These equations ensure that the width, height, and type assigned to each fixture are consistent and that unused fixtures have dimensions and type equal to \(0\).

To respect fixtures availability constraints, we impose:
\[
\sum_{i=1}^{\nmax} T_{i,j} \le \nfix_j, \quad \forall j \in [1..\ntypes] \qquad\qquad
\nmin \le \sum_{i=1}^{\nmax} \used_i \le \nmax
\]

To determine whether two fixtures lie on the same bar for all \(i \in [1..\nmax]\) and \(j \in [i+1..\nmax]\) we impose the following constraints:
\[
\begin{aligned}
	&\samebar_{i,j} \le \used_i \qquad \samebar_{i,j} \le \used_j \\
	&\samebar_{i,j} = \aboveinbar_{i,j} + \belowinbar_{i,j}\\
	&cx_i - cx_j \le M (1 - \samebar_{i,j}) \qquad cx_j - cx_i \le M (1 - \samebar_{i,j}) \\
\end{aligned}
\]

where \(M = \wtab \cdot \htab\) is the \emph{``big-M''} constant. Then to enforce the appropriate safety distances we use the following inequalities:
\[
\begin{aligned}
	&cx_i + \wmin \le cx_j + M(\samebar_{i,j} + \nooverlapright_{i,j}) \\
	&cx_j + \wmin \le cx_i + M(\samebar_{i,j} + \nooverlapleft_{i,j}) \\
	&y_i + h_i + \hmin - y_j \le M (1 - \aboveinbar_{i,j}) \\
	&y_j + h_j + \hmin - y_i \le M (1 - \belowinbar_{i,j}).
\end{aligned}
\]

To prevent overlap between fixtures, for all \(i,j \in [1..\nmax]\) with \(i < j\), we impose:
\[
\begin{aligned}
	&\nooverlapabove_{i,j} + \nooverlapbelow_{i,j} +
	\nooverlapleft_{i,j} + \nooverlapright_{i,j} \ge 1 \\
	&x_i + w_i \le x_j + M (1 - \nooverlapleft_{i,j}) \qquad
	&y_i + h_i \le y_j + M (1 - \nooverlapbelow_{i,j}) \\
	&x_j + w_j \le x_i + M (1 - \nooverlapright_{i,j}) \qquad
	&y_j + h_j \le y_i + M (1 - \nooverlapabove_{i,j})
\end{aligned}
\]
Analogous constraints are introduced to prevent overlap between fixtures and extrusion areas. For each fixture \(i \in [1..\nmax]\) and extrusion area \(j \in [1..\numext]\), the same inequalities are applied using the binary variables \(\aboveext_{i,j}\), \(\belowext_{i,j}\), \(\leftext_{i,j}\), and \(\rightext_{i,j}\).

The workpiece is defined by linear inequalities of the form $\mathtt{a}x + \mathtt{b}y + \mathtt{c} \leq 0$, so for each fixture \(i \in [1..\nmax]\) we impose \(\mathtt{a}\,x'_i + \mathtt{b}\,y'_i + \mathtt{c} \leq 0\), where \(x'_i\) and \(y'_i\) are defined as in the CP model (see Section~\ref{sec:cp_cons}).

To break symmetries, we enforce an ordering on fixture types: $t_i \ge t_{i+1}$ for $i \in [1..\nmax-1]$.

Finally, the absolute value terms appearing in \(\fdist\) are linearized using a big-\(M\) formulation, unfolding each term of the form \(u=|v|\) into: 

$$u \ge v ~\wedge~ u \ge -v ~\wedge~ u \leq v + M (1 - b)  ~\wedge~ u \leq -v + M b~\wedge~ b\in\{0,1\}.$$